\section{Background}
\label{sec:rs-background}

The Recursive Spawning Protocol (RSP) draws from and contributes to five distinct research traditions: agent lifecycle management in multi-agent systems; accountability and provenance in distributed systems; the structural distinction between fixed and mutable delegation chains; hierarchical task decomposition in AI; and the BX3 Framework as a three-layer accountability substrate. This section situates RSP within each tradition, establishes the specific gaps that motivate its design, and demonstrates that these gaps are not addressable by extensions of existing approaches.

% -------------------------------------------------------------------------%
\subsection{Agent Lifecycle Management in Multi-Agent Systems}
% -------------------------------------------------------------------------%

The problem of governing agent populations across their full lifecycle --- from provisioning through decommissioning --- has become a central research concern as autonomous agents move from single-task deployments into complex, multi-agent workflows. The Cloud Native Computing Foundation's 2026 cloud-native agentic standards document identifies \emph{agent identity}, least-privilege access, and tenancy controls as foundational requirements for multi-agent systems operating in Kubernetes environments \cite{cncf-agentic-standards}. The standards specify that agents requiring access beyond a user's permission scope or persisting beyond session termination must carry distinct identities enabling authentication, authorization, and audit enforcement at both identity and execution layers \cite{cncf-agentic-standards}. This identity requirement extends to Just-in-Time (JIT) access provisioning, Attribute-Based Access Control (ABAC), and Policy-Based Access Control (PBAC) for flexible, secure policy definition \cite{cncf-agentic-standards}. Microsoft's Agent Governance Toolkit operationalizes these principles through Agent OS --- a stateless policy engine intercepting agent actions at sub-millisecond latency --- and Agent Mesh, which provides cryptographic decentralized identifiers (DIDs) with Ed25519 signatures, an Inter-Agent Trust Protocol (IATP), and dynamic trust scoring across behavioral tiers \cite{microsoft-agent-governance}. Microsoft's framework demonstrates that structural runtime enforcement of agent identity and trust boundaries is achievable at production scale, but neither the CNCF standards nor Microsoft's toolkit address the structural immutability of the parent delegation chain that RSP provides.

The Agent Lifecycle Protocol (ALP) defines a formal state machine governing six lifecycle stages --- Genesis, Fork, Migration, Retirement, Succession, and Decommission --- with explicit transition rules, boundary hooks, and a genealogical registry \cite{alp2025}. ALP identifies \emph{lineage tracking} and \emph{accountability across generations} as core unsolved challenges, observing that a comprehensive genealogical registry for querying an agent's complete family and history is necessary for audit and compliance but absent from conventional runtimes \cite{alp2025}. RSP directly addresses this gap: the immutable parent pointer is the foundational primitive for the genealogical registry ALP calls for, and the forensic ledger provides the append-only chain of custody that makes lineage verification runtime-computable rather than reconstructed post-hoc. The Human Delegation Protocol (HDP) introduces a JSON-based token binding a human principal to authorization scope and session, with each agent in a multi-hop chain appending signed hops to an append-only provenance record \cite{hdp-draft,hdp-draft-2}. HDP demonstrates that multi-hop, append-only human provenance in agentic workflows is achievable and verifiable offline using only the issuer's public key. RSP adopts HDP's append-only chain principle and extends it structurally: while HDP achieves provenance through signed hop records, RSP achieves structural immutability of the parent pointer at the Fact Layer write protocol level, eliminating the possibility of retroactive chain modification regardless of signing scheme. The core problem RSP addresses is \emph{spawn chain opacity}: the inability to determine, at any point during or after execution, whether a given agent's authority traces to a human principal or to an autonomous predecessor. This opacity arises because parent pointers are mutable in conventional architectures --- they can be redirected by administrative action, policy override, or runtime reconfiguration --- and this mutability is invisible to downstream verifiers who must rely on the chain's integrity without independent confirmation.

% -------------------------------------------------------------------------%
\subsection{Accountability and Provenance in Distributed Systems}
% -------------------------------------------------------------------------%

The accountability problem in distributed systems --- how to construct a tamper-evident record of which actor performed which action, when, and under what authorization --- has been studied extensively in blockchain, database, and operating systems research. The W3C PROV data model formalizes provenance as entities, activities, and agents linked by relations of derivation, attribution, and responsibility \cite{moreau2015-prov}. Under PROV-DM, agents can be attributed to other agents, enabling a provenance graph encoding the full chain of responsibility. The core insight PROV contributes is that the provenance graph is only as reliable as the immutability guarantees of its edges: if the graph is mutable, any attribution derived from it is conditional rather than definitive \cite{moreau2015-prov}.

PROV-AGENT extends the W3C PROV standard with agent-centric entities, activities, and relations, integrating agent tools, prompt/response interactions, model invocations, and telemetry into a unified provenance graph \cite{prov-agent2025}. PROV-AGENT demonstrates convergence between agentic workflow tracking and classical provenance standards, enabling semantic traversals of agent decision chains. However, PROV-AGENT's provenance graph is constructed from runtime observations and logged events; it does not enforce structural immutability of the parent pointer at provisioning time. The provenance record is accurate only to the extent that the logged events are complete and unmodified --- a property that RSP's Fact Layer write protocol makes structurally guaranteed rather than conventionally assumed.

The immutable ledger tradition in distributed systems provides the structural solution. Certificate Transparency demonstrates that append-only logging with hash-chaining can produce tamper-evident records without centralized trust \cite{laurie2013certificate}. Bitcoin's blockchain extends this with consensus-based append-only ledgers demonstrating that sufficiently distributed consensus can produce globally agreed records with no trusted central authority \cite{nakamoto2008bitcoin}. BlockAudit demonstrates a PBFT-based blockchain architecture providing accountability, provenance, and trustworthy auditing for enterprise and distributed settings by leveraging blockchain properties to prevent undetected tampering of logs \cite{blockaudit}. Harpocrates extends this with privacy-preserving zero-knowledge proofs that allow public verification of operation validity while concealing sensitive metadata \cite{harpocrates}. The critical distinction RSP makes is between \emph{runtime-verified accountability} and \emph{retroactive reconstruction}: a forensic ledger records events after they occur and can be used to infer what happened, whereas an immutable delegation chain is verified structurally at every spawn event through the Fact Layer pre-commit hook, confirming the parent pointer's integrity and the scope subset relationship before any child agent is provisioned. The Credential Broker for Agents (CB4A) architecture observes that long-lived API credentials delegated to agents create accountability gaps when agents spawn sub-agents or act across domain boundaries, and that structural response mechanisms are required rather than policy conventions \cite{draft-cb4a}. SentinelAgent introduces a Delegation Chain Calculus (DCC) with seven formal properties and demonstrates that six of seven properties remain robust under adversarial stress while intent preservation degrades against sophisticated paraphrasing attacks \cite{arxiv-sentinel-agent}. DCC's formalization of scope-action pairs and delegation depth provides a useful vocabulary for RSP's own correctness properties, though DCC does not provide the structural parent-pointer immutability that RSP requires.

% -------------------------------------------------------------------------%
\subsection{Fixed vs. Mutable Delegation Chains}
% -------------------------------------------------------------------------%

The structural distinction between fixed and mutable delegation chains is the central design choice determining whether accountability is achievable in a multi-agent system. A mutable delegation chain allows the parent pointer $\alpha(n)$ to be rewritten after provisioning --- by administrative action, policy update, or by a sufficiently privileged actor. A fixed (immutable) delegation chain establishes $\alpha(n)$ once at provisioning and enforces immutability through the Fact Layer write protocol.

Mutable delegation chains are the historical norm. OAuth 2.1's token exchange model enables scoped, time-limited delegation but leaves the delegation graph itself mutable: authorized users can re-delegate, administrators can modify scope, and tokens can be revoked retroactively \cite{UCAN}. UCAN (User Controlled Authorization Networks) introduces self-sovereign authorization with a capability model where delegation chains are encoded as capability graphs verified transitively \cite{UCAN}. UCAN's capability graph is mutable: capabilities can be delegated, transferred, and revoked, with revocation propagating through the graph. This is appropriate for human-controlled authorization networks where the principal remains in the loop, but is structurally insufficient for autonomous agent chains where the human principal may not be reachable at decision time.

Google DeepMind's \emph{Intelligent AI Delegation} framework identifies recursive delegation risk as a critical failure mode in long agent chains, proposing that delegation should transfer not just work but also scoped authority, responsibility, and accountability, with monitoring and trust mechanisms that hold when environments change or agents fail \cite{arxiv-intelligent-delegation}. The framework identifies common multi-agent failure modes including delegation methods relying on ``simple heuristics'' that do not dynamically adapt to unexpected changes or robustly handle failures \cite{arxiv-intelligent-delegation}. The IETF draft on identity management for agentic AI identifies recursive delegation risk --- agents creating sub-agents or delegating tasks introducing complex chains of accountability and control --- as a critical future challenge, with preemptive authorization and careful scoping proposed as mitigations \cite{openid-agentic-ai}. RSP provides the structural complement to these approaches: rather than relying on the correctness of individual authorization decisions, RSP enforces that every spawn event satisfies the scope subset constraint and depth bound as structural invariants, making the delegation chain's integrity independent of the correctness of any individual actor's judgment. The fundamental problem with mutable delegation chains in multi-agent systems is \emph{accountability laundering}: a drifted agent's chain is rewritten to appear authorized, retroactively obscuring the original provenance. RSP eliminates accountability laundering by making the parent pointer immutable regardless of the requestor's identity or privileges: no administrative override exists, and the immutability is enforced at the write protocol level before any modification is executed.

% -------------------------------------------------------------------------%
\subsection{Hierarchical Task Decomposition in AI}
% -------------------------------------------------------------------------%

Hierarchical task decomposition --- the decomposition of complex objectives into subtask hierarchies allocated to specialized agents --- is the dominant paradigm for multi-agent orchestration in contemporary AI systems. The design space spans from classical planning formalisms to LLM-native orchestration frameworks.

Classical HTN (Hierarchical Task Network) planning decomposes a high-level task into a hierarchy of subtasks, where each decomposition is a method specifying how a compound task can be replaced by a partially ordered set of subtasks, terminating when all primitive tasks are executable by the system's primitive action set \cite{ghallab2004htn}. Simon's concept of bounded rationality provides the theoretical foundation: due to cognitive limits, agents cannot evaluate all possible courses of action at once and must decompose complex problems into manageable subproblems \cite{simon}. This insight applies directly to modern LLM-based agents: an agent with a complex goal must decompose it because the full goal space exceeds its representational and computational capacity.

Modern LLM-based agents have revived hierarchical task decomposition under the rubric of dynamic, runtime spawning. AgentOrchestra introduces a hierarchical multi-agent framework combining high-level planning with modular, specialized agents, where a central planner decomposes objectives into sub-goals and delegates to role-specific sub-agents, demonstrating that hierarchical decomposition improves task success rates (GAIA: 83.39\%) over flat-agent baselines \cite{arxiv-agentorchestra}. HALO introduces a three-layer hierarchical framework --- high-level planning agent, mid-level role-design agents, and low-level inference agents --- using Monte Carlo Tree Search to explore agent action spaces, outperforming static role assignment by approximately 14.4\% on average across benchmarks \cite{arxiv-halo}. TDAG introduces dynamic task decomposition where sub-task specifications are adaptively generated based on environmental feedback, demonstrating that dynamic decomposition outperforms static task graphs in noisy environments \cite{arxiv-tdag}. The THREAD system treats LLM generation as a dynamic multi-threaded process where a parent thread can spawn child threads to handle sub-tasks and pause until those children return concise results, enabling hierarchical task decomposition while managing context length \cite{thread2025}.

What connects classical HTN and modern LLM-based decomposition is the spawn tree: a hierarchical structure in which each node represents a subtask and edges represent the decomposition relationship. In both traditions, the tree is typically mutable and implicit --- decomposition decisions are not recorded as first-class authorization artifacts, and child nodes are not distinguished from parent nodes in terms of their accountability status. A classical HTN planner's subtasks are subroutine calls without independent accountability chains; an LLM agent's spawned sub-agents are independent processes, but the chain of authorization is implicit in the agent's context window rather than explicit in an immutable warrant structure. Dynamic hierarchical task decomposition in open-loop agentic systems has been identified as a vector for \emph{autonomous drift}: the progressive expansion of agent behavior beyond authorized scope, precisely because decomposition is unconstrained \cite{autoGPT2023}. Recursive Spawning adopts the hierarchical decomposition paradigm but adds the accountability layer absent from both traditions: each node in the decomposition tree is a first-class accountability entity with an immutable parent pointer, a defined termination condition, and a cryptographically verifiable chain back to a human anchor. Decomposition is not merely a planning technique; it is an authorization event that creates a new accountable entity.

% -------------------------------------------------------------------------%
\subsection{The BX3 Framework as Substrate for Immutable Parent Chains}
% -------------------------------------------------------------------------%

The BX3 Framework provides the minimal sufficient three-layer accountability architecture on which the Recursive Spawning Protocol operates \cite{bx3-framework,bx3framework2026}. BX3 defines three functional layers --- the \textit{Purpose Layer}, the \textit{Bounds Engine}, and the \textit{Fact Layer} --- each defined by the properties it must maintain rather than by the type of actor occupying it. The Purpose Layer carries intent, judgment, and final human accountability; the Bounds Engine carries constrained execution within a defined operating range; the Fact Layer carries deterministic enforcement and forensic auditability. RSP is the second named pillar of the BX3 Framework, complementing Loop Isolation, Spatial Firewall, Sandbox Gate, and the Bailout Protocol \cite{bx3-framework,bansal2019}.

The upstream accountability guarantee of the BX3 Framework states that when any node encounters a state it cannot resolve within its provisioned bounds, accountability escalates recursively upward through the system hierarchy, bypassing all machine actors, until it reaches a human anchor \cite{bx3-framework,bx3framework2026}. This guarantee is not a policy choice; it is a structural property of the three-layer architecture. The Fact Layer creates the immutable parent pointer at spawn time, issues the spawn warrant, and maintains the forensic ledger that records every spawn event. The Fact Layer's determinism is what makes the parent pointer architecturally immutable: once written, the pointer cannot be updated because the Fact Layer does not permit modifications to its authorization records. Any attempt to overwrite $\alpha(n)$ is rejected at the protocol level before execution. This contrasts with ACL-based immutability in conventional systems: administrative compromise in ACL-based systems leads directly to immutability violation, whereas BX3's protocol-level enforcement has no administrative override path \cite{bx3-framework,bansal2019}.

The depth bound $D_{\max}$ is a Purpose-Layer-defined, Bounds-Engine-enforced, Fact-Layer-hardened constraint. The Purpose Layer defines the maximum permissible chain depth as part of the spawn authorization policy. The Bounds Engine evaluates each proposed spawn against this limit, returning \texttt{SPAWN\_REFUSED} if the spawn would exceed it. The Fact Layer's pre-commit hook enforces this limit structurally: any spawn operation that would produce $|C| \geq D_{\max}$ is rejected before execution, regardless of the Bounds Engine's return. The constraint is thus enforced three times --- by the layer that defines it, the layer that evaluates against it, and the layer that structurally blocks violations. Removing any layer causes the corresponding guarantee to collapse: without the Purpose Layer, there is no authoritative human anchor and no source of spawn authorization policy; without the Bounds Engine, depth limits become unenforceable conventions subject to override by operational urgency; without the Fact Layer, immutability of the parent pointer and tamper-evidence of the ledger are promises rather than constraints.

The five BX3 pillars collectively establish the conditions under which the upstream accountability guarantee applies: that BX3 systems fail upward into human consciousness, never downward into autonomous chaos. Loop Isolation, Recursive Spawning, Spatial Firewall, and Sandbox Gate address exception conditions anticipated at design time and resolvable within the machine-only execution graph. RSP specifically addresses the structural conditions that enable autonomous drift --- the mutable parent pointer, the absent scope refinement constraint, and the missing depth bound enforcement --- making the drift triad architecturally impossible through the Fact Layer's write protocol immutability, the Purpose Layer's spawn authorization gate, and the Bounds Engine's depth enforcement.
